\begin{resumo}[Abstract]
 \begin{otherlanguage*}{english}

Machine learning models trained on educational microdata may reproduce and amplify structural inequalities without conventional performance metrics revealing the problem. This work presents AuditEdu-XAI, an algorithmic auditing software tool that detects, quantifies and explains biases in predictive models of performance in the Brazilian National High School Exam (ENEM). Developed under the Design Science Research methodology, the tool adopts a decoupled architecture. A batch pipeline processes the microdata of the 2019, 2020 and 2021 editions, trains XGBoost classifiers and computes, exclusively on the test set, Disparate Impact, Demographic Parity, true and false positive rates (Equal Opportunity and Equalized Odds) and TreeSHAP explanations. Results are persisted as static Parquet files and consumed by a Streamlit interface with response times below 2 seconds in local validation. The interface is publicly hosted free of charge on Streamlit Community Cloud and packaged in a Docker container validated by continuous integration. Applied to models with a ROC AUC of about 0.75, the audit revealed that all non-white groups and candidates with a family income of up to 2.5 minimum wages violate the four-fifths rule in all three editions, a result that holds under bootstrap confidence intervals. It also revealed that the model amplifies the race/colour disparity already present in the data in 14 of 15 group-edition combinations, 13 of them with a confidence interval entirely below zero, and that the absence of a computer at home is the most influential attribute in the predictions, with a per-candidate weight that grew between 2019 and 2021, coinciding in time with the pandemic. To validate the artifact, a synthetic bias injection experiment with a negative control was conducted: when 50\% of the positive labels of mixed-race (\textit{pardo}) candidates were flipped, the group's Disparate Impact dropped to near zero, while the ROC AUC changed, on average, by at most 0.032. This result indicates that the tool detects and locates biases to which conventional performance metrics are largely insensitive.

   \vspace{\onelineskip}

   \noindent
   \textbf{Keywords}: algorithmic fairness; explain\-able artificial intelligence; disparate impact; SHAP; ENEM.
 \end{otherlanguage*}
\end{resumo}
